% !TeX program = lualatex
% !TeX encoding = utf8
% !TeX spellcheck = uk_UA
% !TeX root = ../test.tex
% в преамбулі ElMagConspect.tex

%=========================================================
\chapter{Відносність електричного і магнітного полів}\label{\currfilebase}
\graphicspath{{\currfilebase/Pictures}}
%=========================================================
\epigraph{Відомо, що електродинаміка Максвелла --- в сучасному її вигляді --- у застосуванні до рухомих тіл
	приводить до асиметрії, яка невластива самим явищам.
}{А.~Ейнштейн, \enquote{До електродинаміки рухомих тіл}, 1905}

Явище електромагнітної індукції --- та ланка, що остаточно поєднує електрику й магнетизм у єдину теорію. Ми бачили \ref{ElMagInduction}, що ЕРС індукції виникає у двох, здавалося б, різних ситуаціях: коли провідник рухається в незмінному магнітному полі
(причина --- сила Лоренца) і коли нерухомий провідник перебуває у змінному магнітному полі (причина --- вихрове електричне
поле). Чи справді це \emphx{два різні} явища, чи \emphx{два погляди} на одне й те саме явище, побачені з різних систем відліку?

\begin{figure}[h!]\centering
	\includegraphics[width=1\linewidth]{alice-bob-magnet_colorized.jpg}
	\caption{Аліса та Боб, кожен з яких вважає нерухомим себе, по-різному пояснюють причину виникнення того самого струму в
		провіднику: для Аліси це електричне поле рухомого магніту, для Боба --- сила Лоренца з боку магнітного поля
		нерухомого магніту\label{pic:AliceBob}}
\end{figure}


Щоб це з'ясувати, розгляньмо замкнений провідний контур і магніт, які рухаються один відносно одного рівномірно й
прямолінійно, та двох спостерігачів --- Алісу й Боба (рис.~\ref{pic:AliceBob}). Аліса пов'язана з контуром: для неї контур
нерухомий, а повз нього рухається магніт. Боб пов'язаний з магнітом: для нього магніт нерухомий, а рухається контур. Обидва
реєструють \emphx{той самий} струм, але пояснюють його по-різному. Аліса бачить навколо рухомого магніту електричне поле, яке
й рухає заряди в контурі. Боб електричного поля не бачить --- поле нерухомого магніту для нього суто магнітне, --- і струм він
пояснює силою Лоренца, що діє на заряди рухомого контуру.

Тут ми спираємося на \emphx{принцип відносності}: усі інерціальні системи відліку рівноправні, а фізичний результат ---
показ амперметра в контурі --- залежить лише від відносного руху контуру й магніту, а не від того, хто з них «насправді»
рухається. Слово «інерціальні» суттєве: наведений аргумент працює лише для рівномірного руху. Для довільної індукції
(деформація контуру, змінний струм у котушці) два механізми переходом до іншої системи відліку не зводяться; їх об'єднує
вже правило потоку.

Саме з цієї ситуації починає Ейнштейн свою статтю 1905 року, з якої взято епіграф. Далі в ній сказано:
\begin{tornpage}\small\itshape
	Пригадаємо, наприклад, електродинамічну взаємодію між магнітом та
	провідником зі струмом. Спостережуване явище залежить лише від відносного руху провідника і магніту, тоді як,
	згідно зі звичайним уявленням, два випадки, у яких рухається або одне, або інше з цих тіл, повинні бути строго
	розмежовані. Справді, якщо рухається магніт, а провідник знаходиться в спокої, то навколо магніту виникає
	електричне поле, що має деяку кількість енергії, яке в тих місцях, де знаходяться частини провідника, породжує струм.
	Якщо ж магніт перебуває у спокої, а рухається провідник, навколо магніту немає ніякого електричного поля;
	натомість у провіднику виникає електрорушійна сила, якій самій по собі не відповідає ніяка енергія, але яка --- за умови
	передбачуваної тотожності відносного руху в обох випадках, --- викликає електричні струми тієї ж величини і того ж
	напрямку, що й електричне поле в першому випадку.\\

	% TODO: вказати джерело перекладу
	\hfill А.~Ейнштейн, 1905
\end{tornpage}

Отже, асиметрія --- у теорії, а не в явищі. Явище симетричне щодо обміну «рухається магніт» $\leftrightarrow$ «рухається
провідник», а теорія ці два випадки розрізняє: в одному є електричне поле з енергією, в іншому --- лише ЕРС, якій жодна
енергія не відповідає. Тому запитання ставиться не так, чи електричне поле «насправді» існує. Воно таке: як пов'язані між
собою електричне й магнітне поля, які бачать різні спостерігачі, і що в цьому зв'язку від вибору системи відліку не
залежить?

\begin{Attention}
	Із принципу відносності та з того, що Аліса й Боб описують той самий струм через різні поля, випливає: електричне й
	магнітне поля \emphx{окремо} не є абсолютними --- їхній поділ залежить від системи відліку. Але це не означає, що поле
	є лише «способом опису». Реальним є електромагнітне поле \emphx{як ціле}; від системи відліку залежить тільки його
	розклад на електричну й магнітну частини. Далі ми побачимо й величини, які від системи відліку не залежать. Ця думка
	є провідною лінією розділу.
\end{Attention}

%% --------------------------------------------------------
\section{Нерелятивістські закони перетворення полів}
%% --------------------------------------------------------

%% --------------------------------------------------------
\subsection*{Перетворення електричного поля}
%% --------------------------------------------------------

Нехай заряджена частинка в системі відліку $S$ рухається зі швидкістю $\vect v$ в полях $\Efield$ і $\Bfield$. Тоді на неї діє сила
\begin{equation*}
	\vect F = q\left(\Efield + \frac1c\,\vect v\times\Bfield\right).
\end{equation*}
Перейдімо в систему відліку $S'$, яка рухається зі швидкістю $\vect v$, --- у цій системі частинка \emphx{спочиває}. У системі $S'$ на частинку
діє лише сила з боку електричного поля (магнітне поле на нерухому частинку не діє):
\begin{equation*}
	\vect F' = q\Efield'.
\end{equation*}
У нерелятивістському наближенні сила є інваріантом, тобто $\vect F = \vect F'$. Звідси випливає \emphx{перший закон перетворення}:
\begin{equation}\label{eq:EprimeFromE}
	\tcbhighmath{
		\Efield' = \Efield + \frac1c\,\vect v\times\Bfield.
	}
\end{equation}

\begin{Attention}
	Формула~\eqref{eq:EprimeFromE} показує, що «суто електричне» поле в одній системі відліку стає комбінацією електричного й магнітного полів
	в іншій. Це --- прямий наслідок принципу відносності: жодне з полів не є «абсолютним», обидва --- лише різні прояви єдиного
	електромагнітного поля.
\end{Attention}

Перехід від системи відліку $S'$ до системи $S$ отримується зміною знака швидкості:
\begin{equation}\label{eq:EFromEprime}
	\tcbhighmath{
		\Efield = \Efield' - \frac1c\,\vect v\times\Bfield'.
	}
\end{equation}

%% --------------------------------------------------------
\subsection*{Перетворення магнітного поля}
%% --------------------------------------------------------

Із закону Біо--Савара випливає, що магнітне поле заряду, який рухається зі швидкістю $\vect v$, дорівнює
\begin{equation*}
	\Bfield = \frac1c\,\vect v\times\Efield', \qquad \Efield' = \frac{q\vect r}{r^3}.
\end{equation*}
Розглянемо систему нерухомих (у системі відліку $S'$) заряджених частинок. Вони створюють електростатичне поле
\begin{equation*}
	\Efield' = \sum_k \frac{q_k}{r_k^3}\vect r_k = \sum_k \Efield'_k.
\end{equation*}
Перейдімо в систему відліку $S$, відносно якої система $S'$ рухається зі швидкістю $\vect v$. У $S$ заряди рухаються зі
швидкістю $\vect v$, тому кожен із них створює магнітне поле $\Bfield_k = \frac1c\,\vect v\times\Efield'_k$, а всі разом вони
створюють поле
$\Bfield_k = \frac1c\,\vect v\times\Efield'_k$, а всі разом вони створюють поле
\begin{equation*}
	\Bfield = \frac1c\,\vect v\times\Efield'.
\end{equation*}
Таким чином, у системі відліку, в якій заряди рухаються, виникає магнітне поле. Якщо у вихідній системі присутнє магнітне поле $\Bfield'$
(створюване, наприклад, власним магнітним моментом частинки), то сумарне магнітне поле дається формулою
\begin{equation}\label{eq:BFromBprime}
	\tcbhighmath{
		\Bfield = \Bfield' + \frac1c\,\vect v\times\Efield'.
	}
\end{equation}
Обернений перехід від системи $S$ до системи $S'$ отримується зміною знака швидкості:
\begin{equation}\label{eq:BprimeFromB}
	\tcbhighmath{
		\Bfield' = \Bfield - \frac1c\,\vect v\times\Efield.
	}
\end{equation}

\begin{Attention}
	Пари формул~\eqref{eq:EprimeFromE},~\eqref{eq:BprimeFromB} та~\eqref{eq:EFromEprime},~\eqref{eq:BFromBprime} дають \emphx{нерелятивістські
	закони перетворення} полів при переході від однієї інерціальної системи відліку до іншої. Ці закони справедливі лише при $v \ll c$; точні
	(релятивістські) формули, розглянуті далі, переходять у них у граничному випадку малих швидкостей.
\end{Attention}

\begin{table}[h!]
	\centering
	\caption{Нерелятивістські перетворення полів}
	\label{tab:NonRelFieldTransform}
	\begin{tblr}%
			{
			colspec={X[c,m, bg=yellow!10]X[c,m, bg=yellow!10]},
			row{1}={c, fg=white, bg=themecolorlight, font=\bfseries},
			hlines,
            vlines
			}
		Перехід від $S$ до $S'$ & Перехід від $S'$ до $S$ \\
		$\Efield' = \Efield + \dfrac1c\,\vect v\times\Bfield$ & $\Efield = \Efield' - \dfrac1c\,\vect v\times\Bfield'$ \\
		$\Bfield' = \Bfield - \dfrac1c\,\vect v\times\Efield$ & $\Bfield = \Bfield' + \dfrac1c\,\vect v\times\Efield'$
	\end{tblr}
\end{table}

%% --------------------------------------------------------
\section{Релятивістські закони перетворення полів}
%% --------------------------------------------------------

\subsection*{Інваріантність заряду}

При виведенні закону збереження заряду,
\begin{equation*}
	\frac{\partial\rho}{\partial t} + \divg\vect j = 0,
\end{equation*}
передбачалося, що величина заряду не залежить від його швидкості. Це твердження є фактично результатом досліду. Припустимо, що це не так,
тобто що заряд залежить від швидкості, $q = q(v)$. Візьмімо електронейтральне тіло, в якому додатні й від'ємні заряди точно скомпенсовані.
Нагріватимемо це тіло. Тоді легкі частинки почнуть рухатися швидше, ніж важкі (наприклад, у газі $v_T = \sqrt{3kT/m} \sim m^{-1/2}$). У
результаті заряд одного знака почне переважати --- тіло набуде нескомпенсованого заряду.

Однак такого ефекту не спостерігалося в жодному експерименті. Це означає \emphx{незалежність заряду від швидкості}.

\begin{Attention}
	Інваріантність заряду --- один із фундаментальних принципів електродинаміки. Саме завдяки їй стає можливим узгоджене релятивістське
	перетворення густини заряду й густини струму: заряд елемента об'єму не змінюється, хоча сам об'єм зазнає лоренцевого скорочення.
\end{Attention}

%% --------------------------------------------------------
\subsection*{Перетворення полів}
%% --------------------------------------------------------

Нерелятивістські формули ми вивели, вважаючи силу однаковою в обох системах відліку. Тепер зробимо те саме точно: замінимо
рівність $\vect F = \vect F'$ правильним перетворенням сили. Введемо позначення
$\vect\beta = \vect v/c$, $\gamma = 1/\sqrt{1-(v/c)^2}$ і, як і раніше, вважатимемо, що система $S'$ рухається відносно $S$
зі швидкістю $\vect v$ (вісь $x$ напрямлена вздовж $\vect v$). Індексом «$\parallel$» позначимо компоненти полів, паралельні
вектору швидкості, а індексом «$\perp$» --- перпендикулярні до нього. Виведення спирається на три положення: означення полів
через силу Лоренца, яке однакове в усіх інерціальних системах; інваріантність заряду; принцип відносності, за яким обернене
перетворення отримується заміною $\vect v \to -\vect v$.

\emphx{Перетворення сили.} Розглянемо частинку, яка в даний момент спочиває в системі $S'$. Енергія й імпульс
перетворюються як компоненти 4-вектора $(\varepsilon/c,\ \vect p)$:
\begin{equation*}
	dp_\parallel = \gamma\left(dp'_\parallel + \frac{v}{c^2}\,d\varepsilon'\right), \qquad dp_\perp = dp'_\perp,
	\qquad dt = \gamma\left(dt' + \frac{v}{c^2}\,dx'\right).
\end{equation*}
Для частинки, що спочиває в $S'$, зміщення $dx' = 0$, а $d\varepsilon' = \vect F'\cdot d\vect r' = 0$. Тому
$dt = \gamma\,dt'$, $dp_\parallel = \gamma\,dp'_\parallel$, $dp_\perp = dp'_\perp$, і, ділячи, отримуємо
\begin{equation*}
	F_\parallel = F'_\parallel, \qquad \vect F_\perp = \frac{\vect F'_\perp}{\gamma}.
\end{equation*}

\emphx{Електричне поле.} У системі $S$ частинка рухається зі швидкістю $\vect v$, тому
$\vect F = q(\Efield + \vect\beta\times\Bfield)$; у системі $S'$ вона спочиває, тому $\vect F' = q\Efield'$ (заряд
інваріантний). Вектор $\vect\beta\times\Bfield$ перпендикулярний до $\vect v$, тож
із $F_\parallel = F'_\parallel$ випливає $E_\parallel = E'_\parallel$, а з $\vect F_\perp = \vect F'_\perp/\gamma$ ---
$\Efield_\perp + \vect\beta\times\Bfield = \Efield'_\perp/\gamma$. Отже,
\begin{equation}\label{eq:RelE}
	\Efield'_\parallel = \Efield_\parallel, \qquad \Efield'_\perp = \gamma\left(\Efield_\perp + \vect\beta\times\Bfield\right).
\end{equation}
Як і в розділі~1, точку простору можна вибрати довільно: у ній завжди можна розмістити пробну частинку.

\emphx{Магнітне поле, перпендикулярні компоненти.} Формула~\eqref{eq:RelE} справедлива для будь-яких полів. За принципом
відносності обернене перетворення відрізняється лише знаком швидкості:
$\Efield_\perp = \gamma\left(\Efield'_\perp - \vect\beta\times\Bfield'\right)$. Підставимо в нього $\Efield'_\perp$
з~\eqref{eq:RelE}:
\begin{equation*}
	\Efield_\perp = \gamma^2\left(\Efield_\perp + \vect\beta\times\Bfield\right) - \gamma\,\vect\beta\times\Bfield'
	\quad\Rightarrow\quad
	\vect\beta\times\Bfield' = \gamma\beta^2\,\Efield_\perp + \gamma\,\vect\beta\times\Bfield,
\end{equation*}
де використано $\gamma^2 - 1 = \gamma^2\beta^2$. Помножимо обидві частини векторно на $\vect\beta$ зліва. Для будь-якого
вектора $\vect X$ маємо $\vect\beta\times(\vect\beta\times\vect X) = -\beta^2\vect X_\perp$, тому
$-\beta^2\Bfield'_\perp = \gamma\beta^2\,\vect\beta\times\Efield - \gamma\beta^2\,\Bfield_\perp$, тобто
\begin{equation}\label{eq:RelBperp}
	\Bfield'_\perp = \gamma\left(\Bfield_\perp - \vect\beta\times\Efield\right).
\end{equation}

\emphx{Магнітне поле, паралельна компонента.} Нехай довгий соленоїд, вісь якого напрямлена вздовж $\vect v$,
спочиває в $S'$. Поле всередині нього $B'_\parallel = \frac{4\pi}{c}\,i'$, де $i' = \sigma'u'$ --- лінійна густина струму
(струм на одиницю довжини соленоїда), $\sigma'$ --- поверхнева густина заряду носіїв, $u'$ --- їхня азимутальна швидкість.
У системі $S$ соленоїд рухається вздовж своєї осі. Поперечні розміри не змінюються, поздовжні скорочуються, тому через
інваріантність заряду $\sigma = \gamma\sigma'$. Азимутальна швидкість носіїв, перпендикулярна до $\vect v$, через уповільнення часу
зменшується: $u = u'/\gamma$. Отже, $i = \sigma u = i'$. Соленоїд залишається електронейтральним (струм перпендикулярний до
$\vect v$), тому в $S$ є лише магнітне поле, і $B_\parallel = B'_\parallel$.

Отже, ми отримуємо закони перетворення при переході від $S$ до $S'$:
\begin{equation}\label{eq:RelTransformDirect}
	\tcbhighmath{
		\begin{aligned}
			\Efield'_\parallel &= \Efield_\parallel, \qquad
			\Efield'_\perp = \gamma\left(\Efield_\perp + \vect\beta\times\Bfield\right), \\[4pt]
			\Bfield'_\parallel &= \Bfield_\parallel, \qquad
			\Bfield'_\perp = \gamma\left(\Bfield_\perp - \vect\beta\times\Efield\right).
		\end{aligned}
	}
\end{equation}
Обернені перетворення, від $S'$ до $S$, отримуються заміною $\vect v \to -\vect v$:
\begin{equation}\label{eq:RelTransformParallel}
	\tcbhighmath{
		\begin{aligned}
			\Efield_\parallel &= \Efield'_\parallel, \qquad
			\Efield_\perp = \gamma\left(\Efield'_\perp - \vect\beta\times\Bfield'\right), \\[4pt]
			\Bfield_\parallel &= \Bfield'_\parallel, \qquad
			\Bfield_\perp = \gamma\left(\Bfield'_\perp + \vect\beta\times\Efield'\right).
		\end{aligned}
	}
\end{equation}

\begin{Attention}
	При $v \ll c$ маємо $\gamma \to 1$, і формули~\eqref{eq:RelTransformDirect} та~\eqref{eq:RelTransformParallel} переходять
	у нерелятивістські формули~\eqref{eq:EprimeFromE}--\eqref{eq:BprimeFromB}: паралельні компоненти взагалі не змінюються, а
	для перпендикулярних залишається лише змішування $\vect\beta\times\Bfield$ і $\vect\beta\times\Efield$. Множник $\gamma$
	з'являється лише для компонент, перпендикулярних до руху, --- саме там працюють скорочення довжин і уповільнення часу.
	Формули можна перевірити на прикладі рухомого зарядженого листа далі: він дає $E = \gamma E'$ і $B = \gamma\frac vc E'$,
	у повній згоді з~\eqref{eq:RelTransformParallel}.
\end{Attention}

%% --------------------------------------------------------
\subsection*{Інваріанти електромагнітного поля}
%% --------------------------------------------------------

З формул перетворення~\eqref{eq:RelTransformDirect} випливає, що є два \emphx{інваріанти} електромагнітного поля:
\begin{equation}\label{eq:FieldInvariants}
	\tcbhighmath{
		I_1 \equiv E^2 - B^2 = \text{inv}, \qquad I_2 \equiv \Efield\cdot\Bfield = \text{inv}.
	}
\end{equation}

\emphx{Доведення.} Розкладемо поля на паралельні й перпендикулярні до $\vect v$ складові. Для векторів, перпендикулярних до $\vect\beta$,
скористаємося тотожністю $(\vect\beta\times\vect a)\cdot(\vect\beta\times\vect b) = \beta^2\,\vect a_\perp\cdot\vect b_\perp$, а також тим, що
вектори $\vect\beta\times\Efield$ і $\vect\beta\times\Bfield$ перпендикулярні до $\vect\beta$, а $\vect\beta\times\Efield\perp\Efield$,
$\vect\beta\times\Bfield\perp\Bfield$.

Інваріант $I_2$. Паралельні компоненти не змінюються, тому
\begin{equation*}
	\Efield'\cdot\Bfield' = E_\parallel B_\parallel
	+ \gamma^2\left(\Efield_\perp + \vect\beta\times\Bfield\right)\cdot\left(\Bfield_\perp - \vect\beta\times\Efield\right).
\end{equation*}
Розкриємо добуток. Доданки $\Efield_\perp\cdot(\vect\beta\times\Efield)$ і $(\vect\beta\times\Bfield)\cdot\Bfield_\perp$
дорівнюють нулю, а $(\vect\beta\times\Bfield)\cdot(\vect\beta\times\Efield) = \beta^2\,\Efield_\perp\cdot\Bfield_\perp$. Тому
\begin{equation*}
	\Efield'\cdot\Bfield' = E_\parallel B_\parallel + \gamma^2\left(1-\beta^2\right)\Efield_\perp\cdot\Bfield_\perp
	= E_\parallel B_\parallel + \Efield_\perp\cdot\Bfield_\perp = \Efield\cdot\Bfield,
\end{equation*}
оскільки $\gamma^2(1-\beta^2) = 1$.

Інваріант $I_1$. Аналогічно,
\begin{equation*}
	E'^2 - B'^2 = E_\parallel^2 - B_\parallel^2 + \gamma^2\left[\left(\Efield_\perp + \vect\beta\times\Bfield\right)^2
	- \left(\Bfield_\perp - \vect\beta\times\Efield\right)^2\right].
\end{equation*}
Різниця квадратів у дужках дорівнює
\begin{equation*}
	\left(1-\beta^2\right)\left(E_\perp^2 - B_\perp^2\right)
	+ 2\left[\Efield\cdot(\vect\beta\times\Bfield) + \Bfield\cdot(\vect\beta\times\Efield)\right].
\end{equation*}
Вираз у квадратних дужках дорівнює $\vect\beta\cdot(\Bfield\times\Efield) + \vect\beta\cdot(\Efield\times\Bfield) = 0$, тому
$E'^2 - B'^2 = E_\parallel^2 - B_\parallel^2 + E_\perp^2 - B_\perp^2 = E^2 - B^2$.

\begin{Attention}
	Інваріанти~\eqref{eq:FieldInvariants} показують, що поділ поля на електричну й магнітну частини не абсолютний, але
	деякі його характеристики від системи відліку не залежать: знак $E^2-B^2$ і обертання в нуль добутку $\Efield\cdot\Bfield$
	однакові для всіх спостерігачів. Наступні наслідки випливають саме з цього.
\end{Attention}

%% --------------------------------------------------------
\paragraph*{Наслідки}
%% --------------------------------------------------------

Якщо в якій-небудь системі відліку $I_2 \neq 0$, то не існує системи відліку, в якій одне з полів оберталося б у нуль (тоді
$\Efield\cdot\Bfield$ дорівнювало б нулю). Іншими словами, неможливо вибором системи відліку виключити хоча б одне з полів.

Нехай у системі $S$ інваріант $I_2 = 0$, тобто $\Efield\perp\Bfield$. Якщо при цьому $I_1 > 0$, тобто $|\Efield| > |\Bfield|$,
то існує система відліку $S'$, в якій присутнє тільки електричне поле. Вона рухається відносно $S$ з такою швидкістю, що
\begin{equation*}
	\Bfield' = \Bfield - \frac1c\,\vect v\times\Efield = 0 \quad\Rightarrow\quad \vect v = c\,\frac{\Efield\times\Bfield}{E^2}.
\end{equation*}
Ця швидкість перпендикулярна і до $\Efield$, і до $\Bfield$, тому обидва поля в $S$ поперечні, множник $\gamma$ на умову
$\Bfield' = 0$ не впливає, і формула точна; при цьому $v = cB/E < c$. Величину поля в $S'$ дає інваріант $I_1$: $E' = \sqrt{E^2 - B^2}$.

Якщо ж $I_1 < 0$, то підхожим вибором системи відліку можна виключити електричне поле, так що залишається тільки магнітне
поле. Швидкість відповідної системи відліку така, що
\begin{equation*}
	\Efield' = \Efield + \frac1c\,\vect v\times\Bfield = 0 \quad\Rightarrow\quad \vect v = c\,\frac{\Efield\times\Bfield}{B^2},
\end{equation*}
причому $v = cE/B < c$, а $B' = \sqrt{B^2 - E^2}$.

Нарешті, якщо $I_1 = 0$ і $I_2 = 0$ (наприклад, поле плоскої хвилі), то в будь-якій системі відліку $E = B$, і жодним вибором
системи відліку не можна виключити ні електричне, ні магнітне поле.

% Взаємне розташування векторів $\{\Efield, \Bfield, \vect v\}$ показано на рис.~\ref{pic:FieldsVelocity}.
%
% \begin{figure}[h!]\centering
% %	\localinput{fields-velocity.tikz}
%	\caption{Відносне розташування векторів полів і швидкості системи відліку, в якій одне з полів обертається в нуль
%		(у вихідній системі $\Efield \perp \Bfield$)\label{pic:FieldsVelocity}}
% \end{figure}

%% --------------------------------------------------------
\subsection*{Перетворення потенціалів}
%% --------------------------------------------------------

Електромагнітне поле характеризується скалярним $\phi$ і векторним $\vect A$ потенціалами, причому
\begin{equation*}
	\Efield = -\grad\phi - \frac1c\frac{\partial\vect A}{\partial t}, \qquad \Bfield = \rot\vect A.
\end{equation*}
Ця пара потенціалів повністю визначає електромагнітне поле і, виявляється, утворює \emphx{4-вектор}
\begin{equation*}
	A = (\phi, A_x, A_y, A_z),
\end{equation*}
компоненти якого перетворюються при переході від однієї інерціальної системи відліку до іншої відповідно до перетворень Лоренца:
\begin{equation}\label{eq:PotentialTransform}
	\tcbhighmath{
		\begin{aligned}
			\phi &= \frac{\phi' + (v/c)A'_x}{\sqrt{1-(v/c)^2}}, \qquad
			A_x = \frac{A'_x + (v/c)\phi'}{\sqrt{1-(v/c)^2}}, \qquad
			A_y = A'_y, \qquad A_z = A'_z; \\[8pt]
			\phi' &= \frac{\phi - (v/c)A_x}{\sqrt{1-(v/c)^2}}, \qquad
			A'_x = \frac{A_x - (v/c)\phi}{\sqrt{1-(v/c)^2}}, \qquad
			A'_y = A_y, \qquad A'_z = A_z.
		\end{aligned}
	}
\end{equation}
Передбачається, що вісь $x$ напрямлена вздовж вектора швидкості $\vect v$.

\begin{Attention}
	Формули~\eqref{eq:PotentialTransform} --- це не що інше, як \emphx{перетворення Лоренца} для 4-вектора потенціалу. Саме тому потенціали
	$\phi$ і $\vect A$ нерозривно пов'язані між собою: при переході до рухомої системи відліку скалярний потенціал «перемішується» з
	компонентою векторного потенціалу вздовж напрямку руху. Це ще один прояв єдності електричного й магнітного полів.
\end{Attention}

%\begin{Example}{Магнітне поле рухомого зарядженого листа}
%	Як приклад застосування цих формул розглянемо таку задачу. Нехай тонкий плоский лист діелектрика рівномірно заряджений з поверхневою
%	густиною $\sigma'$ (рис.~\ref{pic:MovingSheet}). Знайдемо магнітне поле, створюване цим листом у тому випадку, коли лист рухається (у
%	своїй площині) зі швидкістю $\vect v$.
%
%	\begin{center}%[1.5][h!]\centering
%%		\localinput{moving-sheet.tikz}
%		\captionof{figure}{Широкий довгий лист діелектрика, рівномірно заряджений по поверхні, рухається рівномірно і прямолінійно зі швидкістю
%			$\vect v$\label{pic:MovingSheet}}
%	\end{center}
%
%	Вважатимемо, що вісь $x$ напрямлена вздовж вектора швидкості листа, а вісь $z$ --- перпендикулярно листу. Величини, що належать до
%	власної системи відліку листа, позначаємо штрихом.
%
%	У власній системі відліку зарядженого листа маємо тільки електричне поле:
%	\begin{equation*}
%		\Efield' = 2\pi\sigma'\vect k,
%	\end{equation*}
%	де $\vect k$ --- одиничний вектор, напрямлений уздовж осі $z$ (розглядаємо тільки область $z > 0$). Потенціал електричного поля в цій
%	системі відліку
%	\begin{equation*}
%		\phi' = -2\pi\sigma' z'.
%	\end{equation*}
%	Векторний потенціал $\vect A' = 0$. У системі відліку, в якій лист рухається, маємо
%	\begin{equation*}
%		\phi = -\frac{2\pi\sigma' z}{\sqrt{1-(v/c)^2}}, \qquad
%		A_x = -\frac{2\pi\sigma' z}{\sqrt{1-(v/c)^2}}\frac{v}{c}.
%	\end{equation*}
%	Ураховано, що при даному переході координата $z$ не змінюється: $z' = z$. Обчислення електричного поля дає:
%	\begin{equation*}
%		\Efield = -\grad\phi = -\vect k\frac{\partial\phi}{\partial z} = \frac{2\pi\sigma'}{\sqrt{1-(v/c)^2}}\vect k.
%	\end{equation*}
%	Маючи на увазі, що векторний потенціал має єдину ненульову компоненту і залежить тільки від координати $z$, знаходимо магнітне поле:
%	\begin{equation*}
%		\Bfield = \nabla\times\vect A =
%		\begin{vmatrix}
%			\vect i & \vect j & \vect k \\
%			0 & 0 & \partial/\partial z \\
%			A_x & 0 & 0
%		\end{vmatrix}
%		= \vect j\frac{\partial A_x}{\partial z}
%		= -\vect j\frac{2\pi\sigma'}{\sqrt{1-(v/c)^2}}\frac{v}{c}.
%	\end{equation*}
%	Таким чином, магнітне поле є однорідним і напрямлене паралельно площині листа і перпендикулярно до швидкості руху, як показано на
%	рис.~\ref{pic:MovingSheet}.
%\end{Example}

%Отримані формули для полів $\Bfield$ і $\Efield$ можна переписати в іншому вигляді, якщо врахувати, що:
%\begin{enumerate}
%	\item лінійні масштаби скорочуються, а відповідна поверхнева густина заряду зростає:
%	\begin{equation*}
%		\sigma = \frac{\sigma'}{\sqrt{1-(v/c)^2}};
%	\end{equation*}
%	\item величина $i = \sigma v$ є \emphx{лінійною густиною струму}, створюваного рухомими зарядами.
%\end{enumerate}
%У результаті знаходимо магнітне поле:
%\begin{equation}\label{eq:MovingSheetB}
%	\tcbhighmath{
%		B_y = -\frac{2\pi}{c}\sigma v = -\frac{2\pi}{c}i.
%	}
%\end{equation}
%Електричне ж поле рухомого листа дається, як і в електростатиці, формулою
%\begin{equation}\label{eq:MovingSheetE}
%	\tcbhighmath{
%		E = 2\pi\sigma,
%	}
%\end{equation}
%але з перетвореною поверхневою густиною заряду $\sigma$.
%
%\begin{Attention}
%	Результат~\eqref{eq:MovingSheetB} має простий фізичний зміст: магнітне поле рухомого зарядженого листа --- це не що інше, як поле
%	нескінченної площини зі струмом, яке ми вивчали в магнітостатиці. Той самий фізичний об'єкт --- заряджений лист --- у різних системах
%	відліку виглядає або як суто електричне джерело, або як електричне й магнітне джерело одночасно. Це ще раз підтверджує відносність
%	поділу поля на електричну та магнітну складові.
%\end{Attention}

%% --------------------------------------------------------
\section{Тензор електромагнітного поля}
%% --------------------------------------------------------

Закони перетворення~\eqref{eq:RelTransformDirect} мають дивний вигляд: шість компонент $\Efield$ і $\Bfield$ переходять одна в
одну за різними правилами для поздовжніх і поперечних складових. Це означає, що вони є компонентами одного об'єкта, як три
просторові й одна часова компоненти 4-вектора. Цей об'єкт --- \emphx{тензор електромагнітного поля}.

\emphx{Позначення.} Координати $x^\mu = (ct,\,x,\,y,\,z)$, метрика з сигнатурою $(+,-,-,-)$, тобто для 4-вектора
$a^\mu = (a^0, \vect a)$ маємо $a_\mu = (a^0, -\vect a)$; грецькі індекси пробігають значення $0,1,2,3$, латинські ---
$1,2,3$, за повторюваними індексами йде сума. Похідна $\partial_\mu = \left(\frac{1}{c}\frac{\partial}{\partial t},\ \nabla\right)$.
4-потенціал: $A^\mu = (\phi, \vect A)$, $A_\mu = (\phi, -\vect A)$.

\emphx{Означення.} Тензор поля будується з потенціалів як
\begin{equation}\label{eq:FieldTensorDef}
	F_{\mu\nu} = \partial_\mu A_\nu - \partial_\nu A_\mu .
\end{equation}
Він не змінюється при калібрувальному перетворенні $A_\mu \to A_\mu + \partial_\mu\chi$. Виразимо його компоненти через
$\Efield$ і $\Bfield$:
\begin{equation*}
	F_{0i} = \partial_0 A_i - \partial_i A_0 = -\frac1c\frac{\partial A^i}{\partial t} - \frac{\partial\phi}{\partial x^i} = E_i,
	\
	F_{12} = \partial_1 A_2 - \partial_2 A_1 = -\left(\frac{\partial A_y}{\partial x} - \frac{\partial A_x}{\partial y}\right) = -B_z,
\end{equation*}
і аналогічно $F_{13} = B_y$, $F_{23} = -B_x$. У вигляді таблиці:
\begin{equation}\label{eq:FieldTensorMatrix}
	\tcbhighmath{
		F_{\mu\nu} =
		\begin{pmatrix}
			0 & E_x & E_y & E_z \\
			-E_x & 0 & -B_z & B_y \\
			-E_y & B_z & 0 & -B_x \\
			-E_z & -B_y & B_x & 0
		\end{pmatrix},
		\qquad
		F^{\mu\nu} =
		\begin{pmatrix}
			0 & -E_x & -E_y & -E_z \\
			E_x & 0 & -B_z & B_y \\
			E_y & B_z & 0 & -B_x \\
			E_z & -B_y & B_x & 0
		\end{pmatrix}.
	}
\end{equation}
Тензор антисиметричний, $F_{\nu\mu} = -F_{\mu\nu}$, тому має рівно шість незалежних компонент --- стільки, скільки в
$\Efield$ і $\Bfield$ разом. Перехід до $F^{\mu\nu}$ означає зміну знака компонент з одним просторовим індексом ($F^{0i} = -F_{0i}$).

\emphx{Перетворення полів.} Компоненти тензора перетворюються за правилом, однаковим для кожного індексу:
\begin{equation*}
	F'^{\mu\nu} = \Lambda^\mu{}_\alpha\,\Lambda^\nu{}_\beta\,F^{\alpha\beta}, \qquad
	\Lambda =
	\begin{pmatrix}
		\gamma & -\gamma\beta & 0 & 0 \\
		-\gamma\beta & \gamma & 0 & 0 \\
		0 & 0 & 1 & 0 \\
		0 & 0 & 0 & 1
	\end{pmatrix},
\end{equation*}
де $\Lambda$ --- перетворення Лоренца при переході до системи $S'$, що рухається вздовж осі $x$ зі швидкістю $v$. Розпишемо
дві компоненти:
\begin{equation*}
	F'^{02} = \gamma F^{02} - \gamma\beta F^{12}
	\ \Rightarrow\ -E'_y = -\gamma E_y + \gamma\beta B_z
	\ \Rightarrow\ E'_y = \gamma\left(E_y - \beta B_z\right),
\end{equation*}
\begin{equation*}
	F'^{12} = -\gamma\beta F^{02} + \gamma F^{12}
	\ \Rightarrow\ -B'_z = \gamma\beta E_y - \gamma B_z
	\ \Rightarrow\ B'_z = \gamma\left(B_z - \beta E_y\right).
\end{equation*}
Це складові формул~\eqref{eq:RelTransformDirect}: $(\vect\beta\times\Bfield)_y = -\beta B_z$, $(\vect\beta\times\Efield)_z = \beta E_y$.
Решта чотири компоненти дають $E'_x = E_x$, $B'_x = B_x$, $E'_z = \gamma(E_z + \beta B_y)$, $B'_y = \gamma(B_y + \beta E_z)$,
тобто \emphx{всі} шість співвідношень~\eqref{eq:RelTransformDirect} одразу.

\begin{Attention}
	Закони перетворення полів, які ми виводили в кілька кроків, є наслідком одного факту: $F^{\mu\nu}$ --- тензор. Причина
	в силі Лоренца: у 4-вимірній формі вона записується як
	\begin{equation*}
		\frac{dp^\mu}{d\tau} = \frac qc\,F^{\mu\nu}u_\nu, \qquad u^\mu = \gamma\left(c,\ \vect v\right),
	\end{equation*}
	де $\tau$ --- власний час. Ліворуч 4-вектор, а $q$ й $u_\nu$ --- інваріант і 4-вектор, тому $F^{\mu\nu}$ мусить
	перетворюватися як тензор. При $\mu = 1,2,3$ праворуч виходить $\gamma q(\Efield + \frac1c\vect v\times\Bfield)$,
	тобто звичайна сила Лоренца, а при $\mu = 0$ --- робота поля $q\,\Efield\cdot\vect v$.
\end{Attention}

%% --------------------------------------------------------
\subsection*{Інваріанти з тензора}
%% --------------------------------------------------------

Так само, як скалярний добуток $a_\mu b^\mu$ двох 4-векторів однаковий в усіх інерціальних системах, скалярами є й повні
згортання тензора зі самим собою. Побудуємо їх.

\emphx{Інваріант $I_1$.} Розглянемо згортання $F_{\mu\nu}F^{\mu\nu}$. Компоненти з індексом $0$ дають
$F_{0i}F^{0i} + F_{i0}F^{i0} = -2E^2$, а чисто просторові --- $F_{ij}F^{ij} = 2B^2$ (кожна компонента $B$ входить двічі).
Тому
\begin{equation}\label{eq:I1Tensor}
	\tcbhighmath{
		F_{\mu\nu}F^{\mu\nu} = 2\left(B^2 - E^2\right)
		\quad\Longleftrightarrow\quad
		I_1 = E^2 - B^2 = -\frac12\,F_{\mu\nu}F^{\mu\nu}.
	}
\end{equation}
Це скаляр, бо $F'_{\mu\nu}F'^{\mu\nu} = F_{\mu\nu}F^{\mu\nu}$: при перетворенні кожен нижній індекс компенсується верхнім, тому що
$\Lambda^{-1}\Lambda = 1$.

\emphx{Інваріант $I_2$.} Другий скаляр потребує \emphx{дуального тензора}
\begin{equation*}
	\tilde F^{\mu\nu} = \frac12\,\epsilon^{\mu\nu\rho\sigma}F_{\rho\sigma},
\end{equation*}
де $\epsilon^{\mu\nu\rho\sigma}$ --- цілком антисиметричний символ Леві-Чівіти, $\epsilon^{0123} = +1$. Обчислення дають, що $\tilde F^{\mu\nu}$
отримується з $F^{\mu\nu}$ заміною $\Efield \to \Bfield$, $\Bfield \to -\Efield$. Зокрема $\tilde F^{01} = F_{23} = -B_x$. Повне
згортання
\begin{equation*}
	\epsilon^{\mu\nu\rho\sigma}F_{\mu\nu}F_{\rho\sigma}
	= 8\left(F_{01}F_{23} - F_{02}F_{13} + F_{03}F_{12}\right)
	= -8\left(E_xB_x + E_yB_y + E_zB_z\right),
\end{equation*}
тому
\begin{equation}\label{eq:I2Tensor}
	\tcbhighmath{
		I_2 = \Efield\cdot\Bfield = -\frac14\,\tilde F^{\mu\nu}F_{\mu\nu} = -\frac18\,\epsilon^{\mu\nu\rho\sigma}F_{\mu\nu}F_{\rho\sigma}.
	}
\end{equation}
Символ $\epsilon^{\mu\nu\rho\sigma}$ не змінює вигляду при власних перетвореннях Лоренца ($\det\Lambda = 1$), тому й цей скаляр
не залежить від системи відліку. Він змінює знак при дзеркальному відбитті (псевдоскаляр): $\Efield$ --- вектор, $\Bfield$
--- аксіальний вектор, і їхній скалярний добуток парність не зберігає. Величина $E^2 - B^2$ від парності не залежить.

Незалежно від дуального тензора: $\det F_{\mu\nu} = (\Efield\cdot\Bfield)^2$. Оскільки $F' = \Lambda F\Lambda^{T}$ і
$\det\Lambda = 1$, визначник не змінюється, і $|I_2|$ інваріантний.

\emphx{Чому їх рівно два.} Складемо комплексний вектор $\vect{\mathcal F} = \Efield + i\Bfield$. Тоді
\begin{equation*}
	\vect{\mathcal F}^{\,2} = E^2 - B^2 + 2i\,\Efield\cdot\Bfield = I_1 + 2i\,I_2,
\end{equation*}
а формули~\eqref{eq:RelTransformDirect} набувають вигляду
$\vect{\mathcal F}'_\parallel = \vect{\mathcal F}_\parallel$,
$\vect{\mathcal F}'_\perp = \gamma\left(\vect{\mathcal F}_\perp - i\,\vect\beta\times\vect{\mathcal F}\right)$.
Це «обертання» комплексного вектора навколо $\vect v$ на уявний кут: $\gamma = \cosh\eta$, $\gamma\beta = \sinh\eta$. Прямою підстановкою
$\vect{\mathcal F}'^{\,2} = \mathcal F_\parallel^2 + \gamma^2(1-\beta^2)\,\vect{\mathcal F}_\perp^{\,2} = \vect{\mathcal F}^{\,2}$
(перехресний член $\vect{\mathcal F}_\perp\cdot(\vect\beta\times\vect{\mathcal F}_\perp)$ дорівнює нулю). Тобто $I_1$ і $I_2$ ---
дійсна й уявна частини єдиного комплексного інваріанта, а обертання комплексного вектора не змінює лише його квадрат.

\begin{Attention}
	Рівняння Максвелла в тензорному вигляді складаються з двох рівнянь:
	\begin{equation*}
		\partial_\mu F^{\mu\nu} = \frac{4\pi}{c}\,j^\nu, \qquad \partial_\mu\tilde F^{\mu\nu} = 0, \qquad j^\nu = (c\rho,\ \vect j).
	\end{equation*}
	При $\nu = 0$ перше з них дає $\divg\Efield = 4\pi\rho$, при $\nu = 1,2,3$ --- $\rot\Bfield - \frac1c\frac{\partial\Efield}{\partial t} =
	\frac{4\pi}{c}\vect j$; друге дає $\divg\Bfield = 0$ і $\rot\Efield + \frac1c\frac{\partial\Bfield}{\partial t} = 0$. У першому
	рівнянні обидві частини є 4-векторами ($\partial_\mu F^{\mu\nu}$ --- згортання 4-вектора $\partial_\mu$ з тензором), а в
	другому 4-вектор дорівнює нулю. Тому рівняння зберігають вигляд у всіх інерціальних системах, якщо $F^{\mu\nu}$ --- тензор, а
	$j^\nu$ --- 4-вектор. Це стисла схема розв'язання задачі~\ref{prob:MaxwellInvariance}.
\end{Attention}

%% --------------------------------------------------------
\section*{Підсумки}
%% --------------------------------------------------------

Усі результати цього розділу спираються на одну спільну ідею: \emphx{електричне й магнітне поля --- це дві грані єдиного електромагнітного
поля}, а їхній поділ залежить від вибору системи відліку. Нерелятивістські закони перетворення~\eqref{eq:EprimeFromE}--\eqref{eq:BprimeFromB}
справедливі при $v \ll c$ і показують, як «суто електричне» поле набуває магнітної складової при переході до рухомої системи відліку.
Релятивістські закони~\eqref{eq:RelTransformParallel} уточнюють ці співвідношення з урахуванням лоренцевого фактора $\sqrt{1-(v/c)^2}$.
Інваріанти~\eqref{eq:FieldInvariants} $I_1 = E^2 - B^2$ та $I_2 = \Efield\cdot\Bfield$ залишаються незмінними в будь-якій інерціальній системі
відліку й визначають, чи можна вибором системи відліку повністю виключити одне з полів. Нарешті, потенціали $\phi$ і $\vect A$ утворюють
4-вектор~\eqref{eq:PotentialTransform}, компоненти якого перетворюються за законами Лоренца, --- це найглибше вираження єдності
електромагнітного поля.

%% --------------------------------------------------------
\section*{Задачі}
%% --------------------------------------------------------

У всіх задачах система $S'$ рухається відносно $S$ зі швидкістю $\vect v$, якщо не сказано інше.

% ЗМІНЕНО: виправлена відповідь
\begin{problem}\label{prob:FieldTransformNonRel}
	\emphx{Нерелятивістське перетворення полів.} У системі відліку $S$ є однорідне електричне поле $\Efield = E_0\vect i$ і магнітне поле
	$\Bfield = B_0\vect j$. Знайдіть поля $\Efield'$ і $\Bfield'$ у системі $S'$, що рухається зі швидкістю $\vect v = v\vect i$ ($v \ll c$).
	\emphx{Відповідь:} $\Efield' = E_0\vect i + \dfrac{v}{c}B_0\vect k$, $\Bfield' = B_0\vect j$.
\end{problem}

\begin{problem}\label{prob:FieldInvariants}
	\emphx{Інваріанти поля.} Доведіть, що величини $I_1 = E^2 - B^2$ та $I_2 = \Efield\cdot\Bfield$ справді є інваріантами відносно
	перетворень~\eqref{eq:RelTransformParallel}. \emphx{Вказівка:} розкладіть поля на паралельну та перпендикулярну до $\vect v$ складові й
	скористайтесь явними формулами перетворення.
\end{problem}

% ЗМІНЕНО: замість виродженої задачі (нульове поле в обох системах) --- потенціали рівномірно рухомого заряду
%\begin{problem}\label{prob:PotentialTransform}
%	\emphx{Потенціали рівномірно рухомого заряду.} Точковий заряд $q$ спочиває на початку координат системи $S'$, яка рухається
%	відносно $S$ зі швидкістю $\vect v = v\vect i$; початки координат і відліку часу збігаються при $t = t' = 0$. У системі $S'$
%	потенціали $\phi' = q/r'$, $\vect A' = 0$, де $r' = \sqrt{x'^2 + y'^2 + z'^2}$.
%	\begin{enumerate}
%		\item Скориставшись перетвореннями~\eqref{eq:PotentialTransform} і співвідношеннями
%		$x' = (x - vt)/\sqrt{1-(v/c)^2}$, $y' = y$, $z' = z$, знайдіть $\phi$ і $\vect A$ у системі $S$.
%		\item Обчисліть $\Efield$ і $\Bfield$ та покажіть, що $\Bfield = \frac1c\,\vect v\times\Efield$.
%		\item Переконайтесь, що при $v \ll c$ поле $\Bfield$ збігається з полем точкового заряду, яке дає закон Біо--Савара.
%	\end{enumerate}
%
%	\emphx{Відповідь:} $\phi = \dfrac{q}{\sqrt{(x-vt)^2 + \bigl(1-(v/c)^2\bigr)(y^2+z^2)}}$, $\vect A = \dfrac vc\,\phi\,\vect i$
%	(тобто $A_y = A_z = 0$); $\Bfield = \frac1c\,\vect v\times\Efield$.
%\end{problem}

\begin{problem}\label{prob:PureEandB}
	\emphx{Чисте електричне поле в одній системі --- чисто магнітне в іншій?} Чи може так статися, що в деякій інерціальній системі відліку
	існує \emphx{тільки} електричне поле ($\Efield \neq 0$, $\Bfield = 0$), а в іншій інерціальній системі --- \emphx{тільки} магнітне поле
	($\Efield' = 0$, $\Bfield' \neq 0$)? Обґрунтуйте відповідь, використовуючи інваріанти поля.
\end{problem}

% ЗМІНЕНО: уточнено, відносно чого визначається швидкість
\begin{problem}\label{prob:PureElectricExist}
	\emphx{Коли існує система з чисто електричним полем?} Нехай у деякій інерціальній системі відліку задано однорідні поля $\Efield$ і
	$\Bfield$. За яких умов на ці поля існує така інерціальна система, в якій магнітне поле дорівнює нулю? Знайдіть її швидкість
	відносно вихідної системи.

	\emphx{Відповідь:} система з $\Bfield' = 0$ існує тоді й лише тоді, коли $\Efield \perp \Bfield$ та $E > B$. Швидкість такої системи
	\begin{equation*}
		\vect v = c\,\frac{\Efield\times\Bfield}{E^2}.
	\end{equation*}
\end{problem}

\begin{problem}\label{prob:NeutralWireParadox}
	\emphx{«Нейтральний» провід зі струмом.} У системі відліку $S$ довгий прямий провід нерухомий і електрично нейтральний: додатні іони
	спочивають, а електрони дрейфують зі швидкістю $u$, створюючи струм $I$. Покажіть, що в системі відліку $S'$, яка рухається вздовж
	проводу зі швидкістю $v$ (меншою за $u$), провід виявляється \emphx{зарядженим}. Як це узгоджується з інваріантністю заряду?

	\emphx{Вказівка:} скористайтесь релятивістським перетворенням густини заряду й густини струму як компонент 4-вектора $j = (\rho c, \vect j)$
	та лоренцевим скороченням відстаней між іонами й електронами. Ця задача --- класична ілюстрація того, що магнітна сила в одній системі
	відліку є проявом електричної сили в іншій.
\end{problem}

\begin{problem}\label{prob:MaxwellInvariance}
	\emphx{Чи зберігають рівняння Максвелла форму?} Покажіть, що рівняння Максвелла
	\begin{equation*}
		\divg\Efield = 4\pi\rho, \qquad \rot\Bfield - \frac1c\frac{\partial\Efield}{\partial t} = \frac{4\pi}{c}\vect j,
	\end{equation*}
	\begin{equation*}
		\divg\Bfield = 0, \qquad \rot\Efield + \frac1c\frac{\partial\Bfield}{\partial t} = 0
	\end{equation*}
	зберігають свій вигляд при переході до будь-якої інерціальної системи відліку, якщо поля перетворюються
	за формулами~\eqref{eq:RelTransformParallel}, а густини заряду й струму --- за перетвореннями Лоренца для 4-вектора $j = (\rho c, \vect j)$.
	\emphx{Вказівка:} це найпряміший шлях довести, що рівняння Максвелла \emphx{релятивістськи інваріантні} --- і що саме тому вони
	«автоматично» містять у собі спеціальну теорію відносності, задовго до її створення.
\end{problem}

\begin{Questions}
\begin{quizlist}
	\item У чому полягає парадокс Аліси й Боба про магніт і провідник? Чому в обох випадках виникає той самий струм, попри різне фізичне пояснення (сила Лоренца проти вихрового електричного поля)?
	\item Що мав на увазі Ейнштейн, кажучи, що асиметрія --- <<в теорії, а не в явищі>>? Чому питання про те, чи електричне поле <<насправді>> існує, поставлене некоректно?
	\item Виведіть нерелятивістський закон перетворення електричного поля $\Efield' = \Efield + \dfrac1c\vect v\times\Bfield$, виходячи з інваріантності сили в нерелятивістському наближенні.
	\item Виведіть нерелятивістський закон перетворення магнітного поля $\Bfield = \Bfield' + \dfrac1c\vect v\times\Efield'$, спираючись на закон Біо--Савара для рухомих зарядів.
	\item Чому інваріантність заряду (незалежність від швидкості) --- фундаментальний експериментальний факт? Як довід із нагріванням електронейтрального тіла демонструє, що було б, якби заряд залежав від швидкості?
	\item На яких трьох положеннях ґрунтується релятивістське виведення законів перетворення полів? Виведіть перетворення сили $F_\parallel=F'_\parallel$, $\vect F_\perp = \vect F'_\perp/\gamma$ для частинки, що спочиває в $S'$.
	\item Виведіть релятивістський закон перетворення паралельної та перпендикулярної складових електричного поля, $\Efield'_\parallel=\Efield_\parallel$, $\Efield'_\perp=\gamma(\Efield_\perp+\vect\beta\times\Bfield)$.
	\item Виведіть релятивістське перетворення перпендикулярної складової магнітного поля $\Bfield'_\perp = \gamma(\Bfield_\perp - \vect\beta\times\Efield)$, користуючись принципом відносності й уже знайденим перетворенням для $\Efield$.
	\item На прикладі рухомого соленоїда покажіть, чому паралельна складова магнітного поля не змінюється при переході в іншу інерціальну систему відліку, $B_\parallel = B'_\parallel$.
	\item Запишіть повну систему релятивістських перетворень полів. Чому при $v\ll c$ вони переходять у нерелятивістські формули, і чому множник $\gamma$ з'являється лише для перпендикулярних компонент?
	\item Що таке інваріанти електромагнітного поля $I_1=E^2-B^2$ та $I_2=\Efield\cdot\Bfield$? Доведіть інваріантність $I_2$, використовуючи розклад полів на паралельну й перпендикулярну складові.
	\item За яких умов на інваріанти $I_1$, $I_2$ можна вибором системи відліку повністю виключити електричне поле, а за яких --- магнітне? Виведіть відповідну швидкість такої системи відліку.
	\item Чому для поля плоскої хвилі ($I_1=I_2=0$) жодним вибором системи відліку не можна виключити ні електричне, ні магнітне поле?
	\item Покажіть, що скалярний $\phi$ і векторний $\vect A$ потенціали утворюють 4-вектор, компоненти якого перетворюються за формулами Лоренца.
	\item Дайте означення тензора електромагнітного поля $F_{\mu\nu} = \partial_\mu A_\nu - \partial_\nu A_\mu$ через 4-потенціал. Чому цей тензор антисиметричний і має рівно шість незалежних компонент?
	\item Виведіть явний вигляд матриці $F_{\mu\nu}$ через компоненти $\Efield$ і $\Bfield$. Покажіть, як з перетворення тензора Лоренца відновлюються всі шість формул релятивістського перетворення полів.
	\item Як інваріанти $I_1$ і $I_2$ виражаються через повні згортання тензора поля $F_{\mu\nu}F^{\mu\nu}$ та дуального тензора $\tilde F^{\mu\nu}F_{\mu\nu}$? Чому $I_2$ --- псевдоскаляр, а $I_1$ --- звичайний скаляр?
	\item Складіть комплексний вектор $\vect{\mathcal F} = \Efield + i\Bfield$ і поясніть, чому саме тому інваріантів рівно два --- дійсна й уявна частини квадрата цього вектора.
	\item Як записуються рівняння Максвелла в коваріантному тензорному вигляді $\partial_\mu F^{\mu\nu} = (4\pi/c)j^\nu$ та $\partial_\mu\tilde F^{\mu\nu}=0$? Чому з цього запису випливає релятивістська інваріантність рівнянь Максвелла?
	\item Розгляньте парадокс <<нейтрального>> провідника зі струмом: чому в системі відліку, що рухається вздовж проводу, він виявляється зарядженим, і як це узгоджується з інваріантністю заряду?
\end{quizlist}
\end{Questions}